\documentclass[11pt]{article}
\usepackage[letterpaper,margin=.82in,top=.85in,bottom=.8in]{geometry}
\usepackage{graphicx,booktabs,xcolor,hyperref,fancyhdr,lmodern}
\definecolor{ink}{HTML}{19374B}\definecolor{teal}{HTML}{116C67}
\hypersetup{colorlinks=true,urlcolor=teal}
\pagestyle{fancy}\fancyhf{}\lhead{\small UIC MATH 496 \quad / \quad MURL Chess}\rhead{\small Fall 2026}
\cfoot{\small Akbar Aman \quad / \quad Research contribution \hfill \thepage}
\setlength{\headheight}{14pt}\setlength{\parindent}{0pt}\setlength{\parskip}{6pt}
\newcommand{\heading}[1]{{\color{ink}\Large\bfseries #1}\par}
\newcommand{\tagline}[1]{{\color{teal}\small\bfseries #1}\par}
\newcommand{\sub}[1]{\par\smallskip{\color{teal}\bfseries #1}\par}
\newcommand{\fig}[1]{\includegraphics[width=\linewidth]{../deliverables/november05/figures/#1.png}}
\newcommand{\repo}{https://github.com/ak23bar/murl-chess}
\begin{document}
\tagline{RESEARCH FINDINGS \quad / \quad NOVEMBER 2026}
{\color{ink}\Huge\bfseries Finding the first move in a forced mate}\par
{\large Runtime, accuracy, and failure modes of Minimax and Monte Carlo tree search}\par
Akbar Aman \quad $\cdot$ \quad Prepared for La Min Aung, group poster coordinator\\
November 3, 2026 \quad $\cdot$ \quad Group findings due November 5
\sub{Research question}
On the same 20 mate-in-three positions, how long do two search algorithms take, how often do they choose the unique winning first move, and what can an actual mistake tell us about the limits of each search?
\begin{center}\begin{tabular}{lrrrr}\toprule
October setting & Calls & Finished & Finished key hits & Median time$^*$\\\midrule
Minimax, 3 plies &20&20&6&0.535 s\\
Minimax, 5 plies &20&2&2&8.195 s\\
MCTS, 1,000 simulations &60&60&9&2.724 s\\\bottomrule
\end{tabular}\end{center}
{\small $^*$Among finished calls. MCTS used three seeds on each of the same 20 positions. The five-ply median describes only puzzles 16 and 18.}
\sub{What the measurements support}
Three-ply Minimax returned the puzzle key on six positions. MCTS returned it in nine trials across five distinct positions. In puzzle 6 all three seeds chose to take a rook with the queen. Each missed a queen check that forces mate. This contrast is visible on the board and in the recorded root visits.

Five plies reach the attacker's third move. When the full search finished, it found both keys. The 20-second ceiling stopped the other 18 calls. Eight partial answers matched the key, but the work needed to finish those searches remains unknown.

A separate November run gave the algorithms the same time ceilings. Increasing the ceiling from a quarter second to one second increased the number of returned keys from 4 to 5 across 20 Minimax decisions and from 3 to 6 across 60 MCTS decisions. These observations describe this puzzle set, implementation, machine, and run order.

\href{\repo}{Repository, original reports, machine-readable tables, and poster figures}
\newpage
\tagline{01 \quad / \quad STUDY DESIGN}
\heading{How the moves were scored}
\sub{From course diagrams to checkable positions}
The September study measured runtime on eight deliberately varied boards without a known optimal move for each. October used the corrected course sheet, \textit{Positions2.pdf}. Its 20 diagrams supply a different experiment with a specific objective. All 1,280 squares were checked against the diagram glyphs. Position 3 has Black to play. The others have White to play. FEN records use the indicated side, no castling or en passant rights, and clocks 0 and 1. The October sensitivity check for positions 1 and 4 preserved each unique key under the three tested alternative Black castling-rights settings.

A \emph{ply} is one move by one player. The attacker moves on plies 1, 3, and 5. At attacker nodes the verifier needs one winning continuation. At defender nodes every legal reply must still lose within the horizon. The exact verifier checked every legal first move at bounds 1, 3, and 5. Each position had no shorter forced mate and exactly one winning first move at five plies. A separate certificate checker replays a winning strategy against every legal defense. Both use the same chess rules library as the search programs.

The November audit also replayed all five plies of every printed main line. All 20 keys agree with the verifier, and all 20 lines end in checkmate. Thirteen printed lines differ from the stored sample continuation. A sample shows one path through a strategy. The certificate covers the defender's alternatives.
\sub{Two ways of deciding}
Plain Minimax searches every legal move to its depth, maximizing for White and minimizing for Black. Checkmate takes precedence over material. At a nonterminal leaf it counts pawns, minor pieces, rooks, and queens as 1, 3, 5, and 9. Equal evaluations retain the first legal move. On timeout it keeps the best fully evaluated root branch, or the first legal move if none finished.

MCTS expands a random untried move, uses UCT to revisit branches, and simulates random legal play for up to 40 plies. Terminal outcomes yield a win, loss, or draw reward. At the cutoff it converts material balance to a bounded reward. It returns the most visited root move, breaking ties by the mover's reward. A sampled favorable outcome can occur even when a defender has a reply that defeats the puzzle objective.
\sub{What counts as accuracy}
A hit means the returned move equals the unique verified key. It does not establish that a search found the complete continuation. October distinguishes completed key hits from returned keys on timeout. November asks a different operational question: which move was available when time ran out?
\newpage
\tagline{02 \quad / \quad EXECUTION TIME AND ACCURACY}
\heading{Search effort depends on the position}
\fig{runtime}
{\small Figure 1. October wall time by puzzle. Each square gives the median across three MCTS seeds. Triangles mark the 20-second ceiling. The eventual full-depth times for those positions are censored.}
\sub{RQ1. Execution time}
Three-ply Minimax finished all 20 searches. The MCTS median across the 60 individual calls was 2.724 s versus 0.535 s for Minimax. The workloads differ: one traverses a fixed depth, the other runs 1,000 simulations. Their medians do not compare equal work. The earlier eight-board study illustrates why piece count alone is inadequate. At depth three, two 16-piece boards had 15 versus 53 legal first moves, visited 2,978 versus 123,342 nodes, and took 0.087 versus 3.586 s at the median.
\sub{RQ2. Move accuracy}
The October completed-call hit rates were 30\% for three-ply Minimax and 15\% for MCTS. Each MCTS seed found three keys, but they were not always the same three. Only five distinct puzzles had a hit at any seed. Treating the 60 trials as 60 different boards overstates the breadth of this evidence.
\sub{Shared time limits, measured November 3}
The run shuffled all 160 calls with order seed 496. Minimax requested three plies. MCTS requested 100,000 iterations so the clock would stop it; seeds were 0, 1, and 2. Every method could return early. The table counts returned decisions including partial calls.
\begin{center}\small\begin{tabular}{lrrrrr}\toprule
Setting &Returned hits &Finished &Hits when finished & Median s & Median sims\\\midrule
0.25 s Minimax &4/20 &4 &1 &0.2501 & --\\
0.25 s MCTS &3/60 &0 &0 &0.2502 &90\\
1.00 s Minimax &5/20 &15 &5 &0.5286 & --\\
1.00 s MCTS &6/60 &0 &0 &1.0002 &341\\\bottomrule
\end{tabular}\end{center}
The shared ceiling gives each algorithm the same opportunity to spend time. Actual time used can differ. Small overshoots arise from cooperative clock checks. The cutoff-dependent moves can change on replay.
\newpage
\tagline{03 \quad / \quad INCORRECT MOVES}
\heading{A defense explains what a miss costs}
\fig{decisions}
{\small Figure 2. Teal arrows show the verified keys. Orange arrows show moves actually returned in October. All three boards have White to play.}
\sub{6. A rook capture misses the back-rank mate}
MCTS chose \textbf{Qxe6} at all three seeds. After \textbf{...Rd8}, White cannot force mate in the remaining three plies. The key is \textbf{Qc8+}: after ...Rd8, Qxd8+ Re8, Qxe8 is mate. At seed 0 the capture received 88 visits versus 56 for the key; mean White rewards were 0.544 and 0.443. Of 1,000 rollouts, 881 reached the material cutoff. These observations explain the final visit-based selection. They do not isolate how often a different rollout policy would succeed.
\sub{13. Taking a knight spends the key tempo}
Three-ply Minimax and all three MCTS seeds played \textbf{Kxd1}. Black can answer \textbf{...Kg8}, escaping the prescribed mate horizon. White instead has \textbf{Nf6}: one printed continuation is ...Ne7, Rxe7 Ne3, Rh7 mate. The five-ply Minimax call returned Nf6 at timeout. It did not finish evaluating the root.
\sub{16. Looking two plies further changes the decision}
Three-ply Minimax chose \textbf{Kg4}; \textbf{...Kg8} leaves no forced mate within three more plies. Five-ply Minimax finished in 10.012 s and chose \textbf{Rf7}. After ...Kg8, Rd7 g5, b8=Q is mate. Here changing depth with the same evaluator exposed the mate and changed the first move. Position 18 finished at both depths, selecting the same key each time.
\sub{A correction to the solution commentary}
For position 18, after Rxg7 Kf8 Ne6+ Ke8, the booklet gives c8=R mate. Its accompanying explanation says queen promotion stalemates. Replay on the printed board shows \textbf{both c8=R and c8=Q give checkmate}; bishop and knight promotion each give stalemate. The first move and all reported accuracy counts are unaffected.
\newpage
\tagline{04 \quad / \quad INTERPRETATION AND REUSE}
\heading{What the poster can conclude}
\sub{Evidence for the third research question}
Depth mattered in puzzle 16, where the complete five-ply search changed a wrong move into a verified key. Complexity mattered for feasibility: position 18 has three legal first moves and took 0.025 s at depth three. Position 4 has 43 and took 3.513 s. The root count measures only the first layer; branching, checks, and terminal positions deeper in the tree also determine work. The evidence does not isolate an independent effect of root mobility on accuracy.

The November ceilings changed search effort. Median MCTS simulations rose from 90 to 341 and more returned moves matched the key. The October logs separately show that 89.56\% of 60,000 rollouts ended at the evaluation cutoff. A heuristic that favors material and a policy that samples random defense are plausible contributors to missed tactical lines. These runs did not independently vary those two mechanisms.
\sub{Limits and attribution}
These are twenty selected tactical compositions with unique keys, nineteen with White to move. They do not estimate tournament strength or general-game accuracy. Five-ply completion was highly selective. There were three seeds per MCTS condition and one timed November run per seed and ceiling. Machine load affects timeouts. All searches use the same Python move generator, so the exact solver is algorithmically separate but cannot independently detect a common rules-engine defect.

This is Akbar Aman's research contribution for the collaborative UIC MURL Chess poster coordinated by La Min Aung. Other members' experiments can be added under their own names, hardware, runs, and definitions of success. The Spring 2026 poster is a layout reference only.
\sub{Reproduction and sources}
From the \href{\repo}{repository}, follow the README to the versioned code, FENs, all 100 October search calls, 60 verifier calls, strategy certificates, and 160 November decisions. Environment records identify Windows 11 build 26200, AMD64 Family 25 Model 80, CPython 3.12.14, and chess 1.11.2. Original September and October PDFs are in \texttt{Reports/}. The corrected diagram sheet and solution booklet are archived with their provenance.

Algorithm background: Russell and Norvig, \textit{Artificial Intelligence: A Modern Approach}, 4th ed., ch.~5; Kocsis and Szepesv\'ari, \href{https://doi.org/10.1007/11871842_29}{Bandit Based Monte-Carlo Planning} (2006). The repository bibliography and source code identify the implementation choices. Numerical claims above come from the course benchmark logs.
\end{document}
